\section{Research Design, Objectives, and Hypotheses}\label{sec:design}
\subsection{Objectives}
The research questions translate into seven objectives.
\begin{enumerate}
\item[O1:] Reproduce DIS-FDM's method and its reported performance trends in an independent Contiki-NG/Cooja implementation, to validate the environment before extending it.
\item[O2:] Implement an attack model in which the same node executes multicast DIS flooding and false-parent DAO manipulation concurrently.
\item[O3:] Extend DIS-FDM's evidence fusion with a self-calibrated DAO-consistency source, calibrate the DIS threshold on each node, and derive the parent-change thresholds from attack-free runs instead of fixing them by hand.
\item[O4:] Evaluate the extension against DIS-FDM, Secure-RPL, RPL-MRC, and Li-MSD in detection, service, and cost, including robustness to attack rate and reception loss.
\item[O5:] Validate every comparison statistically across independent paired seeds in static and mobile networks.
\item[O6:] Give each node a tabular Q-learning decision layer, specified and analyzed mathematically before simulation.
\item[O7:] Define the Attack Resilience Index and the Downward Recovery Time, and evaluate all defenses on them.
\end{enumerate}

Table~\ref{tab:chain} links each gap to the questions, objectives, hypotheses, and contributions that address it, so that every result in Section~\ref{sec:results} can be traced back to the gap it answers.

\begin{table}[t]
\caption{Gap, research question, objective, hypothesis, and contribution chain (contributions numbered as in Section~\ref{sec:introduction}).}\label{tab:chain}
\centering\footnotesize
\setlength{\tabcolsep}{4pt}
\begin{tabularx}{\columnwidth}{@{}lllX@{}}
\toprule
Gap & RQ & Objectives & Hypothesis and contribution \\
\midrule
3 (no accuracy evidence) & RQ1 & O1, O5 & H1; reproduction with TPR, FPR, and tests (C4, C5) \\
1 (single vector) & RQ2 & O2, O7 & H2; combined threat model, ARI and DRT (C3, C4) \\
1, 2 (hand-set evidence) & RQ3 & O3, O4, O7 & H3; calibrated DAO evidence and repair (C1, C4) \\
4 (no on-node learning) & RQ4 & O6 & H4; Q-learning layer and analysis (C1, C2) \\
3 & RQ5 & O4, O5 & H5; robustness across conditions (C4) \\
\bottomrule
\end{tabularx}
\end{table}

\subsection{Hypotheses and Their Development}
Each hypothesis, its measurement, and its decision rule were fixed in the research proposal before the main experiments, and they are decided here exactly as stated there. We give each as a null ($H_0$) and alternative ($H_1$) pair and explain why it was expected. Where the evaluation adds a measurement that the proposal did not tie to a hypothesis (ARI in H2 and H3, attack rate and reception loss in H5), it is reported next to the hypothesis as additional evidence and does not change the decision.

\emph{H1 (reproduction).} DIS-FDM reduced DIO overhead by suppressing the DIO responses that floods trigger \cite{prajisha2025}. An implementation of the same rules in the same simulator family should reproduce the direction of these effects, although absolute values depend on unreported settings. $H1_0$: the reproduced effects differ in direction from the published ones. $H1_1$: they have the same direction and the defenses rank in the same order.

\emph{H2 (combined damage).} DIS flooding consumes control-plane capacity, and false-parent DAOs break downward routes \cite{rajasekar2022,limsd2024}. Because the two mechanisms act on different routing functions, a node that runs both should impose both costs. $H2_0$: the combined attack is no worse than each single attack. $H2_1$: it is significantly worse than each single attack in downward delivery, DIO overhead, radio power, or delay. We also report ARI for H2, since O7 introduces it as our service measure.

\emph{H3 (DAO evidence and repair).} DIS-FDM observes no DAO field, so it cannot see, or undo, a false parent \cite{prajisha2025}. Adding a DAO-consistency source should reveal the manipulation, and repairing the advertised parent should restore the route. $H3_0$: AEF-RPL's TPR and precision do not differ from DIS-FDM's. $H3_1$: AEF-RPL has a significantly higher TPR or precision, with no significant rise in FPR. RQ3 also asks about resilience, so we ran the same tests on downward delivery and ARI and report them under H3.

\emph{H4 (on-node learning).} Under the fixed rule, every node reacts in the same way to the DIS bursts and parent changes produced by a moving neighbor. A Q-learner instead updates its action from the outcomes it sees locally \cite{watkins1992,qsecrpl2023}. We therefore expected learning to isolate fewer honest nodes, mainly in the mobile networks. $H4_0$: learning does not change FPR. $H4_1$: learning lowers FPR without significantly lowering ARI.

\emph{H5 (consistency).} If the gains come from the protocol mechanism and not from a particular topology, they should appear across seeds, attacker fractions, mobility, and network sizes. $H5_0$: the H3 and H4 effects appear only in some conditions. $H5_1$: they are significant in both static and mobile networks and keep their direction at 500 and 1000 nodes. We could not test the 500- and 1000-node part (Section~\ref{sec:limitations}). At 30 nodes, we varied the attack rate and the reception loss instead.

For H2 the expected direction is known (the combined attack should be worse), so its tests are one-sided. H3 and H4 are tested two-sided, since an added mechanism might also degrade detection.

\subsection{Conceptual Framework}
Fig.~\ref{fig:framework} summarizes how the experiments are organized. We treat the defense configuration as the independent variable and change it along three axes: the evidence sources used, the decision layer, and the way DAOs are handled. The dependent variables are grouped into detection, service, and cost, and four moderators set how hard the attack is. Within each paired comparison, seed, layout, traffic, duration, and radio parameters stay fixed. Because of this pairing, an ablation that changes only one axis attributes any difference in the outcome to that axis.

\begin{figure}[t]
\centering
\begin{tikzpicture}[font=\scriptsize,
  box/.style={draw=blue!60!black,rounded corners=2pt,fill=blue!4,align=left,inner sep=3pt,text width=2.35cm},
  mod/.style={draw=orange!70!black,rounded corners=2pt,fill=orange!6,align=left,inner sep=3pt,text width=6.6cm},
  arr/.style={-{Stealth[length=2mm]},thick}]
\node[box] (iv) {\textbf{Independent: defense}\\ Evidence sources (DIS, parent, energy, DAO)\\ Decision layer (fixed, Q-learning)\\ DAO handling (drop, repair)};
\node[box,right=1.9cm of iv,text width=2.6cm] (dv) {\textbf{Dependent}\\ Detection: TPR, FPR, precision, latency\\ Service: downward PDR, ARI, DRT\\ Cost: DIO, radio power, CPU time};
\node[mod,above=0.45cm of $(iv.north)!0.5!(dv.north)$] (mo) {\textbf{Moderators:} attacker fraction (10, 25, 50\%), mobility (0, 30\%), attack interval (3 to 30\,s), reception success (1.0, .95, .8)};
\node[mod,draw=gray,fill=gray!6,below=0.45cm of $(iv.south)!0.5!(dv.south)$] (co) {\textbf{Controls:} paired seed (layout, attackers, mobility), traffic, 60\,min duration, onset at 300\,s, radio ranges, MAC};
\draw[arr] (iv) -- node[pos=.27,above]{H2--H4} (dv);
\coordinate (hit) at ($(iv.east)!0.72!(dv.west)$);
\draw[arr,orange!70!black] (mo.south -| hit) -- node[right]{H5} ($(hit)+(0,0.08)$);
\draw[arr,gray] (co.north -| hit) -- ($(hit)-(0,0.08)$);
\end{tikzpicture}
\caption{Conceptual framework: the defense configuration (independent variable) affects detection, service, and cost (dependent variables) under four moderators, with paired controls.}\label{fig:framework}
\end{figure}
